\section{Related research and the scope of the contribution}
\subsection{Historical expectations, sentiment, and uncertainty}
Historical forecast interpretation has direct precedents. \citet{mathy2017} quantify contemporaneous and forward-looking assessments in the business press during the Great Depression. Their object is close to ours: what economic observers believed while events unfolded. Our longer annual corpus and automated candidate generation extend the scale of this inquiry; they do not originate the use of historical press statements to study expectations. The contribution here is a documented chain from extraction choices to claim interpretation and baseline-relative outcomes.

The broader text-as-data literature provides the statistical setting \citep{gentzkow2019}. Newspaper-derived constructs need not be explicit forecasts. \citet{bybee2024} extract topical attention from business news and examine its relationship with economic activity, including incremental information beyond conventional indicators. \citet{barbaglia2022} use aspect-specific news sentiment for macroeconomic forecasting. These approaches motivate a demanding comparison with macro-only predictors, but their targets differ from the correctness of an attributable statement. A topic can be informative without containing a forecast; an explicit forecast can be incorrect while revealing a valuable historical belief. No common-data replication against those published signals is attempted here.

Financial-news research supplies another relevant distinction. \citet{tetlock2007} links pessimistic language in a Wall Street Journal column to market behavior. \citet{garcia2013} studies news sentiment and stock-return predictability, finding that the predictive relationship is concentrated in recessions. These studies establish substantive precedents for extracting economically informative language. They do not imply that every negative sentence predicts a price decline, or that a sentiment coefficient measures a newspaper forecaster's accuracy. For our application, the relevant comparison must align the quoted proposition, asset or macroeconomic target, forecast origin, and horizon. A sentiment index and a directional claim may contain related information without being interchangeable measurements.

\citet{baker2016} develop newspaper-based measures of economic policy uncertainty and validate their construction using several forms of evidence, including human readings. Their historical US index reaches back to 1900 \citep{bakerhistorical}. It offers a natural comparison for a long-run newspaper corpus, but measures policy-related uncertainty rather than the correctness of particular forecasts. The documented historical index also uses source-specific normalization and overlapping newspaper sets. That design highlights a limitation of our exports: without a complete article-processing denominator, raw forecast counts cannot be interpreted as a coverage-normalized attention index. We do not estimate a relationship with EPU in this paper; adding its series to a future comparison would require a documented alignment and an explicit account of changing source coverage.

\subsection{Professional forecasts and the meaning of confidence}
The Livingston Survey \citep{philalivingston} and Survey of Professional Forecasters (SPF; \citealp{philaspf}) supply established records of explicitly elicited economic expectations. The former began in 1946 and the latter in 1968. Federal Reserve staff projections are separately available in the Philadelphia Fed's Tealbook/Greenbook datasets \citep{philagreenbook}. These records are potential external comparators, rather than observations included in the newspaper sample or benchmarks already tested here. A valid comparison would harmonize the economic variable, units, target period, release vintage, and information date before evaluating relative performance. Survey response dates, newspaper publication dates, and staff forecast dates need not describe identical information sets. In particular, the staff projections become public only after a five-year lag; their historical forecast date alone does not make their contents an admissible contemporary public predictor.

\citet{campbell2024} examine probabilistic SPF forecasts and find evidence of overprecision, with little evidence of optimistic bias. Their analysis distinguishes confidence in outcomes from the direction of forecast errors and uses an exploration--validation design. This distinction matters for our interpretation of linguistic confidence. An extracted label such as confident or hedged is not an elicited probability distribution. A low directional hit rate therefore does not, by itself, demonstrate overprecision, and confident language cannot be assigned a numerical probability without a separately validated mapping. Conversely, the probability produced by our claim-correctness model concerns whether a machine-scored label will be correct; it is not the probability originally expressed by a historical speaker. We make no claim to reproduce the SPF overprecision result with newspaper labels.

\subsection{Text measurement and extraction validation}
Validation must match the intended construct. \citet{shapiro2020} compare sentiment measures with human ratings and separately investigate economic applications. \citet{grimmer2013} emphasize problem-specific validation of automated text analysis. \citet{loughran2011} demonstrate that general-purpose negative-word lists can misclassify language in financial filings, motivating domain-specific measurement. Their evidence concerns dictionaries and 10-Ks, not our newspaper extractor; its relevance is the need to validate the interpretation assigned to language in its own setting. For the present study, agreement about the text and success at predicting an economic series are separate requirements. A model's association with future production cannot substitute for checking that the source actually made the forecast attributed to it.

Low-cost language-model annotation is itself established research. \citet{gilardi2023} compare ChatGPT classifications with human annotations across several text tasks and document cost advantages in those applications. Our LOC comparison evaluates extraction spans and inclusion decisions, while the ProQuest review evaluates historical economic fields. Neither sample establishes that the same model ranking transfers across domains. The contribution is therefore application-specific evidence, including adverse results, rather than a claim that inexpensive automated annotation is newly possible.

\subsection{Forecast evaluation, revisions, and model selection}
Proper scoring rules provide a foundation for evaluating probabilistic forecasts \citep{gneiting2007}. Brier loss and log loss assess probabilities, whereas directional accuracy evaluates categorical decisions and AUC evaluates ranking. These measures answer different questions. A proper probability score cannot repair an incorrectly defined outcome: the economic target must first correspond to the proposition being evaluated. Likewise, differences in target prevalence make scores from different samples unsuitable as direct measures of a model's improvement. Our comparisons therefore emphasize matched observations, explicit targets, and the same baseline predictions under each stated rule.

Historical prediction also requires careful treatment of data revisions. \citet{croushore1999} explain how real-time vintages differ from later datasets. We distinguish print-time predictor information, labels admitted to model fitting, and revised outcomes used for retrospective assessment. The selected historical macro benchmark reconstructs vintage information for fitting; the claim-correctness exercise uses frozen revised assessment series and a release allowance. These are different evidentiary standards, and we do not describe the latter as a complete recreation of historical real-time forecasting.

Finally, \citet{cawley2010} show how optimization of a finite-sample model-selection criterion can introduce selection bias in reported performance. Chronological separation addresses information ordering, but cannot restore an untouched evaluation sample after its results have influenced analysis choices. This motivates our disclosure of prior outcome exposure and the exploratory status of the saved extension. The concern is selection and interpretation, not a general objection to simulation or to evaluating estimators against known population parameters. Appendix~\ref{app:examples} states that distinction explicitly.

